\scanpage{10}
\mspage{7}
\uline{avec pour courbe universelle} (\uline{en coordonnées non homogènes}) \handwritten{\tikz[baseline=-.2ex,x=1ex,y=1ex]{\draw (0,0)--(.32,1.18) .. controls (.15,1.27) and (.46,1.55) .. (.54,1.38) .. controls (.62,1.27) and (.40,1.13) .. (.32,1.18);}}\inserted{\handwritten{section $e$}}
\[
Y^2=4X^3-g_2X-g_3
\]
\uline{et pour différentielle invariante} $\pi=d\,X/Y$.

De plus, lorsque 2 et 3 sont inversibles, la donnée d'une forme différentielle invariante suffit à rigidifier une courbe de genre 1.

Il résulte aussitôt de (2.2) qu'on a
\[
c_4^3-c_6^2=1728\,\Delta\tag*{(2.6)}
\]
et cette formule, se réduisant à une identité algébrique, est valable en toute caractéristique.

3.\qquad \uline{Cas où} 2 \uline{et} $c_4$ \uline{sont inversibles}.

Supposons que $E$ soit une courbe de genre un sur une base $S$ où 2 soit inversible. La forme différentielle invariante $\pi$ étant donnée, on peut choisir $x$ et $y$ tels que dans (1.3) on ait $a_1=a_3=0$. Les changements de coordonnées respectant cette condition s'obtiennent en faisant $u=1$, $s=t=0$ dans (1.2).

On a donc
\[
y^2=x^3+a_2x^2+a_4x+a_6\tag*{(3.1)}
\]
\[
b_2=4a_2\qquad b_4=2a_4\qquad b_6=4a_6\qquad d_6=8(a_2a_4-9a_6)
\]
\[
c_4=2^4(a_2^2-3a_4)\qquad -c_6=2^6a_2^3-2^5\,3^2\,a_2a_4+2^5\,3^3\,a_6\quad\text{et}\tag*{(3.2)}
\]
\[
\Delta=2^4\,\delta \qquad ,
\]
% archive sheet number: 9
